\documentclass[conference, 11pt]{IEEEtran}
% Fuente: Arial o Times New Roman 11pt, Interlineado 1.15, Márgenes 2.5cm (aprox. controlado por IEEEtran, se ajusta abajo)
\usepackage[utf8]{inputenc}
\usepackage[spanish]{babel}
\usepackage{graphicx}
\usepackage{amsmath}
\usepackage{booktabs}
\usepackage{url}
\usepackage{float}
\usepackage{setspace}
\usepackage[margin=2.5cm]{geometry}
\setstretch{1.15} % Interlineado 1.15

\begin{document}

\title{Clasificación Médica Multimodal con Mixture of Experts \\ y Ablation Study del Router}

\author{\IEEEauthorblockN{1\textsuperscript{st} Nombre Apellido}
\IEEEauthorblockA{\textit{Ingeniería de Datos / IA} \\
\textit{Universidad}\\
Ciudad, País \\
email@universidad.edu}
\and
\IEEEauthorblockN{2\textsuperscript{nd} Nombre Apellido}
\IEEEauthorblockA{\textit{Ingeniería de Datos / IA} \\
\textit{Universidad}\\
Ciudad, País \\
email@universidad.edu}
}

\maketitle

\begin{abstract}
[Escribe aquí el resumen del proyecto. 1-2 párrafos. Describe brevemente el problema de clasificar imágenes médicas heterogéneas sin metadatos, la arquitectura del sistema Moisture of Experts (MoE) implementada, y el resultado del ablation study compitiendo el ViT contra los métodos estadísticos.]
\end{abstract}

\begin{IEEEkeywords}
Mixture of Experts, Vision Transformer, Deep Learning, Medical Imaging, Ablation Study.
\end{IEEEkeywords}

\section{Introducción}
% 0.5 páginas
[Escribe aquí la Motivación del problema. ¿Por qué es importante unificar imágenes 2D y 3D en un mismo sistema? 
Plantea explícitamente la pregunta científica central: ¿Justifica el Vision Transformer su costo computacional como router frente a métodos estadísticos clásicos operando sobre los mismos embeddings?
Finaliza mencionando tu contribución al proyecto y los experimentos realizados. ¡Ojo! No repitas el enunciado del proyecto.]

\section{Arquitectura Propuesta}
% 1 página
[Describe cómo funciona tu Sistema MoE completo de principio a fin.] \\

\textbf{A. Preprocesador Adaptativo} \\
El sistema implementa un preprocesador dinámico que detecta automáticamente la dimensionalidad del tensor de entrada (imágenes 2D tradicionales vs. volúmenes 3D) omitiendo el uso de metadatos clínicos explícitos. Previo a esta inferencia en tiempo real, los datasets independientes superan etapas exhaustivas de Análisis Exploratorio de Datos (EDA) y preprocesamiento de caja negra.
    
Particularmente para el experto encargado de la evaluación ósea (Osteoartritis de rodilla), las radiografías 2D originales de la Osteoarthritis Initiative presentaban alta variabilidad en resolución, contraste y formato de codificación. El preprocesador unificó estas condiciones exportando los tensores al estándar clínico \texttt{.mha} y estandarizando la resolución iterativa a $224 \times 224$ píxeles. Adicionalmente, dado el bajo contraste intrínseco de los tejidos óseos frente a los tejidos blandos, se aplicó a todo el corpus de entrenamiento un filtro CLAHE (\textit{Contrast-Limited Adaptive Histogram Equalization}) global que mejoró la cuantificación de descriptores críticos como entropía y \textit{edge count}. Esta mejora fue esencial para la posterior diferenciación por red convolucional de los osteofitos y del estrechamiento del espacio articular requeridos por el grado empírico de Kellgren-Lawrence (KL). Paralelamente, en operaciones volumétricas, los flujos son redimensionados espacialmente a dimensiones de $64 \times 64 \times 64$ para que el \textit{router} generalice de manera unificada. \\

\textbf{B. Backbone Compartido} \\
[Describe la arquitectura de tu ViT congelado (o similar) y cómo se extraen los CLS tokens $z \in \mathbb{R}^{d_{model}}$ para alimentar los routers a partir de las imágenes.] \\

\textbf{C. Arquitecturas de los Expertos} \\
Cada experto dentro del sistema \textit{MoE} está diseñado y pre-entrenado para extraer representaciones óptimas de dominios anatómicos específicos. Para el análisis de radiografías de rodilla (Osteoartritis), se seleccionó la arquitectura \textbf{VGG-16 con Batch Normalization (VGG-16 BN)}. Aunque modelos más modernos o más profundos ofrecen menor carga de parámetros, la matriz tradicional de VGG-16 demostró empíricamente una alta sensibilidad en la extracción de características óseas. La incorporación de capas de \textit{Batch Normalization} estabilizó abruptamente la convergencia sobre el dataset de la Osteoarthritis Initiative (OAI), aislando efectivamente las 5 clases ordinales (KL 0 a 4). [Justifica las demás arquitecturas elegidas para NIH ChestX-ray, Páncreas, etc.]

% \textbf{Instrucción Figura 1 Obligatoria}: Diagrama de arquitectura
\begin{figure}[H]
    \centering
    % \includegraphics[width=0.48\textwidth]{figuras/arquitectura_moe.png}
    \vspace{4cm} % Quitar este vspace cuando pongas la imagen
    \caption{Diagrama de arquitectura del sistema MoE implementado (puede ser el mismo diagrama que usas para el dashboard, mostrando la entrada $\rightarrow$ preprocesador $\rightarrow$ backbone ViT $\rightarrow$ router $\rightarrow$ 5 expertos).}
    \label{fig:architecture}
\end{figure}

\section{Ablation Study del Router}
% 1.5 páginas
[Aquí se documenta el experimento principal y científico del proyecto. Explica cómo la entrada al router es siempre la misma (el CLS token extraído en la Fase 0) y cómo comparaste matemáticamente 4 cabezas diferentes de decisión:]
\begin{itemize}
    \item \textbf{ViT + Linear (Softmax)}: Baseline de Deep Learning optimizado por GD.
    \item \textbf{ViT + GMM}: Modelado paramétrico evaluado por EM.
    \item \textbf{ViT + Naive Bayes}: Asunción de independencia paramétrica.
    \item \textbf{ViT + k-NN}: Acercamiento no paramétrico FAISS.
\end{itemize}

[Discusión: ¿Qué método ganó y por qué? Habla de la latencia observada, el consumo de VRAM, y evalúa críticamente tu hipótesis inicial frente a los resultados obtenidos.]

% \textbf{Instrucción Figura 2 Obligatoria}: Tabla comparativa del ablation study
\begin{table}[htbp]
\caption{Tabla Comparativa del Ablation Study en Routers}
\begin{center}
\begin{tabular}{|l|c|c|c|}
\hline
\textbf{Router} & \textbf{\textit{Routing Acc.}} & \textbf{\textit{Latencia}} & \textbf{\textit{VRAM (\textasciitilde)}} \\
\hline
ViT + Linear & 0.00 \% & X ms & X MB \\
ViT + GMM & 0.00 \% & Y ms & Y MB \\
ViT + Naive Bayes & 0.00 \% & Z ms & Z MB \\
ViT + k-NN & 0.00 \% & W ms & W MB \\
\hline
\end{tabular}
\label{tab:ablation}
\end{center}
\end{table}

\section{Resultados de Clasificación}
% 1.5 páginas
El rendimiento individual de cada especialista es fundamental para garantizar que el \textit{Mixture of Experts} se nutra de predicciones robustas. En el caso del experto óseo (VGG-16 BN para Osteoartritis), el entrenamiento se realizó en dos fases para evitar sobreajuste catastrófico: una fase inicial de \textit{feature extraction} seguida de una de \textit{fine-tuning}. Dada la naturaleza desbalanceada del dataset OAI, se utilizó focal loss ponderado inversamente. 

El modelo reportó un \textbf{Accuracy global de 82\%} y un \textbf{F1 Macro-Score de 0.8176} (con un AUC-ROC multiclase de 0.9615). Al analizar el comportamiento clase a clase, el sistema demostró ser altamente competente en casos definitivos: sujetos completamente saludables (KL-0, F1 0.86) y osteoartritis severa (KL-4, F1 0.90). La mayor varianza y ruido se localizó en las clasificaciones adyacentes a KL-2, reflejando la complejidad clínica del estrechamiento óseo incipiente. Al aplicar Grad-CAM y Score-CAM sobre las radiografías 2D, se verificó mediante Inteligencia Artificial Explicable (XAI) que las regiones focales de activación se centraban sobre las roturas cartilaginosas y no en el fondo o ruido exterior, comprobando la eficacia del preprocesamiento CLAHE.

[Detalla los resultados individuales del F1 Macro Score para los demás datasets. Compara cómo reaccionó tu sistema general contra métricas "baseline" de la literatura.]

% \textbf{Instrucción Figura 3 Obligatoria}: Curvas de entrenamiento
\begin{figure}[H]
    \centering
    % \includegraphics[width=0.48\textwidth]{figuras/curvas_entrenamiento.png}
    \vspace{4cm} % Quitar este vspace
    \caption{Curvas de entrenamiento (loss en el eje primario y F1 Macro en el eje secundario) del sistema MoE completo a través de las épocas de fine-tuning gloabl.}
    \label{fig:training_curves}
\end{figure}

[Inserta también la matriz de confusión del experto que obtuvo el menor rendimiento. Analiza el comportamiento, por ejemplo, los falsos positivos y falsos negativos más comunes y a qué crees que se deben.]

\section{Balance de Carga}
% 0.5 páginas
[Documenta cómo implementaste la Auxiliary Loss de Switch Transformer. Explica con qué valor calibraste $\alpha \in [0.01, 0.1]$ y cómo llegaste a ese número.]

% \textbf{Instrucción Figura 4 Obligatoria}: Evolución de balance de carga
\begin{figure}[H]
    \centering
    % \includegraphics[width=0.48\textwidth]{figuras/load_balance.png}
    \vspace{4cm} % Quitar este vspace
    \caption{Evolución del balance de carga de los 5 expertos. Histograma de transcurso de épocas para observar cómo $f_i$ y la fracción $P_i$ convergen, asegurando que $max(f_i)/min(f_i) < 1.30$.}
    \label{fig:load_balance}
\end{figure}

\section{Alcances y Limitaciones}
% 1 página
[Contesta las siguientes inquietudes con un abordaje crítico/científico:] \\
1. \textbf{¿Qué funciona bien y por qué?} \\
2. \textbf{¿Qué NO funcionó como se esperaba?} Discute si tu redimensión 3D causó pérdida de detalle, si FAISS era inestable o la varianza de BatchSize causó ruido con Gradient Accumulation.\\
3. \textbf{Limitaciones Computacionales}: Describir cuellos de botella para cargar tantos datos simultáneos, OOM, y transferencias CPU $\rightarrow$ GPU.\\
4. \textbf{Limitaciones de Imbalance}: Comentar sobre la severidad de tener solo cientos de NIfTI de Páncreas frente a +112K Rx Tórax y las heurísticas como resample.\\

% \textbf{Instrucción Figura 5 Obligatoria}: Attention Heatmap del router
\begin{figure}[H]
    \centering
    % \includegraphics[width=0.48\textwidth]{figuras/attention_heatmaps.png}
    \vspace{4cm} % Quitar este vspace
    \caption{Attention Heatmap superpuesto sobre al menos 3 imágenes de diferentes modalidades para entender visualmente dónde enfocó el CLS Token el ViT Router antes de derivar la imagen.}
    \label{fig:heatmaps}
\end{figure}


\section{Conclusiones}
% 0.5 páginas
[El gran final: Da la respuesta directa y cruda a la pregunta científica que planteaste en la introducción. Añade lecciones prácticas que aprendiste diseñando sistemas Switch Transformers y Mixture of Experts en imágenes médicas volumétricas integradas.]


\begin{thebibliography}{00}
% Debes tener al menos 5 referencias. Se recomienda citar a los creadores del dataset, paper base y paper Switch Transformer de Google
\bibitem{b1} Fedus, W., Zoph, B., \& Shazeer, N. (2021). ``Switch Transformers: Scaling to Trillion Parameter Models.'' \textit{JMLR}. 
\bibitem{b2} Jacobs, R. A., Jordan, M. I., Nowlan, S. J., \& Hinton, G. E. (1991). ``Adaptive Mixtures of Local Experts.'' \textit{Neural Computation}.
\bibitem{b3} Dosovitskiy, A. et al. (2020). ``An Image is Worth 16x16 Words: Transformers for Image Recognition at Scale.'' \textit{ICLR 2021}.
\bibitem{b4} Hihn, H., \& Braun, D. A. (2024). ``Mixture of Experts for Image Classification: What's the Sweet Spot?'' \textit{arXiv:2411.18322}.
\bibitem{b5} Pancreatic Cancer Dataset Creators (Año). Título. [Añadir referencia de Zenodo o LUNA16/NIH según dataset].
\end{thebibliography}

\end{document}
